\begin{textsample}{POS Dim 1 – 2015\_c6 – Score 2.00 – 2015\_c6/enade\_2015\_tec\_gastr\_1938c9be4e.txt}  \label{ex:f1_pos_152}
CAPÍTULO 3 ANÁLISE TÉCNICA DA PROVA Este capítulo tem por objetivo apresentar o desempenho dos estudantes concluintes de Tecnologia em Gastronomia no Enade/2015.
Para isso, foram calculadas as estatísticas básicas da prova em seu todo ( seção 3.1.1 ), bem como as estatísticas relacionadas aos Componentes de Formação Geral ( seção 3.1.2 ) e de Conhecimento Específico da Área ( seção 3.1.3 ).
Dadas as suas características, foram analisadas, em separado, as questões objetivas ( seção 3.2 ) e as questões discursivas ( seção 3.3 ).
Para as questões objetivas, foram disponibilizados os índices de facilidade e de discriminação bisserial, também em separado, para os Componentes de Formação Geral ( seção 3.2.1 ) e de Conhecimento Específico ( seção 3.2.2 ).
De cada componente, uma das questões foi escolhida para exemplificar a análise gráfica, relacionando as alternativas escolhidas pelos estudantes ( inclusive o gabarito ) com o número de acertos no componente.
O Anexo I apresenta a íntegra da análise gráfica para todas as questões objetivas.
Para cada uma das questões discursivas, os conteúdos dos tipos mais comuns de respostas dos estudantes são apresentados e comparados com o padrão de respostas esperado ( ver Anexo VIII com o padrão de respostas ).
Tomando como base as duas questões discursivas do Componente de Formação Geral, a seção 3.3.1.6 apresenta comentários sobre a correção das respostas com respeito à Língua Portuguesa.
Nas Tabelas 3.1 e 3.3 são apresentados o tamanho da população inscrita, ausente, presente e percentual de faltosos, respectivamente, por Grande Região e por Categoria Administrativa e Organização Acadêmica da IES.
As demais tabelas apresentam as seguintes estatísticas das notas¹⁵ : média do desempenho na prova, erro padrão da média, desvio padrão, nota mínima, mediana e nota máxima para cada um de seus componentes.
Tais estatísticas contemplam o total de estudantes concluintes da Área de Tecnologia em Gastronomia inscritos e presentes à prova do Enade/2015, tendo em vista agregações ou por Grandes Regiões e o país como um todo, ou por Categoria Administrativa e Organização Acadêmica da IES.
Em relação aos gráficos de distribuição de notas, o intervalo considerado foi de 10 unidades, aberto à esquerda e fechado à direita, com exceção do primeiro intervalo, [ 0 ; 10 ], fechado em ambos os extremos.
Para os gráficos de distribuição das notas das ¹⁵ Uma definição dessas estatísticas pode ser encontrada no Glossário. questão em branco16 e nota zero. 3.1 CONDIÇÃO DE PRESENÇA E ESTATÍSTICAS BÁSICAS DA PROVA Esta seção apresenta a condição de presença, além das estatísticas selecionadas e histogramas da nota geral ( 3.1.1 ) e de cada componente : Formação Geral ( 3.1.2 ) e Conhecimento Específico ( 3.1.3 ).
São também apresentadas estatísticas selecionadas de subpopulações caracterizadas por Grande Região, Categoria Administrativa e Organização Acadêmica. 3.1.1 Condição de Presença e Estatísticas Básicas Gerais A Tabela 3.1 apresenta a Condição de Presença na prova, por Grande Região, dos estudantes concluintes de Tecnologia em Gastronomia.
A população total de Inscritos foi de 5.717.
Destes, 4.540 estiveram Presentes, sendo 20,6 \% o índice de não comparecimento.
A Região de maior abstenção foi a Centro-Oeste ( 32,7 \%, ) e a de menor abstenção foi a Norte ( 8,5 \% ).
Já a Tabela 3.2 apresenta as Estatísticas Básicas da prova, por Grande Região, dos estudantes concluintes de Tecnologia em Gastronomia.
A Média das notas da prova como um todo ( nas seções seguintes serão analisados os Componentes de Formação Geral e de Conhecimento Específico ) foi 51,5, sendo que os alunos da região Norte obtiveram a Média mais baixa ( 43,5 ), e os das regiões Sudeste e Sul obtiveram a Média mais alta ( 53,2 ).
As demais médias foram : 48,8 na região Nordeste ; e 52,0 na região Centro-Oeste.
O Desvio padrão para o Brasil como um todo foi 13,6, sendo o maior Desvio padrão encontrado na região Norte ( 13,8 ) e o menor, na região Sul ( 13,0 ), indicando uma dispersão um pouco menor das notas desta última região.
A região que obteve a maior nota Máxima foi a Sudeste ( 91,0 ), ao passo que a região que atingiu a menor nota Máxima foi a Norte ( 75,0 ).
A Mediana do Brasil como um todo foi 51,7, sendo a maior Mediana obtida nas regiões Sudeste e Sul ( 53,9 ), e a menor obtida no Norte ( 43,9 ).
A nota Mínima foi zero nas regiões Sudeste e Sul, enquanto a maior nota Mínima foi obtida na região Nordeste ( 8,4 ).
Considerando -se as notas segundo Grande Região, observa -se que existe diferença estatisticamente significativa ao nível de 95 \% entre a menor Média, obtida na região Norte ( 43,5 ), e as médias de todas as demais regiões.
O comportamento das notas dos estudantes de todo o Brasil pode ser observado no Gráfico 3.1 que apresenta um histograma com a distribuição das mesmas.
Essa é uma distribuição unimodal com a moda no intervalo ( 50 ; 60 ].
A Tabela 3.3 apresenta informações referentes à Condição de Presença desagregadas por Categoria Administrativa e por Organização Acadêmica.
Da população total de Inscritos, 5.454 são de IES Privada e 263 de IES Públicas.
Em relação à Organização Acadêmica, a maior participação foi obtida por estudantes de Universidades ( 2.128 ), seguido de Faculdades ( 1.810 ) e Centros Universitários ( 1.779 ).
A Categoria Administrativa de maior abstenção foi a Pública ( 21,7 \% ), e entre as Organizações Acadêmicas foi a das Universidades ( 23,2 \% ), que ficou acima da abstenção nacional de 20,6 \%.
Já a Tabela 3.4 apresenta informações referentes à Média da nota final desagregadas por Categoria Administrativa e por Organização Acadêmica.
A Média das notas da prova como um todo foi 51,5.
Em relação à Categoria Administrativa, os estudantes das IES Públicas obtiveram Média mais alta ( 55,7 ), e os das IES Privadas obtiveram Média mais baixa ( 51,3 ), que a Média nacional.
Observa -se que existe diferença estatisticamente significativa entre as médias das notas das IES Públicas e Privadas.
A diferença entre as médias das regiões Sudeste e Sul, e Norte ( 9,7 ), as maiores e a menor Média, é superior à diferença entre IES Públicas e Privadas ( 4,4 ), caracterizando uma maior diversidade regional do que administrativa.
No tocante à Organização Acadêmica, apenas as Faculdades obtiveram Média mais baixa ( 49,0 ) que a nacional.
As demais médias foram : 53,4 nas Universidades e 51,9 nos Centros Universitários.
O Desvio padrão para as IES Privadas ( 13,6 ) foi igual e para os Centros Universitários ( 13,7 ) foi superior ao do Brasil como um todo ( 13,6 ), indicando uma dispersão um pouco maior das notas nesta Organização Acadêmica.
Constata -se que existe diferença estatisticamente significativa ao nível de 95 \% nas médias das notas dos estudantes das Faculdades em relação aos demais tipos de Organização Acadêmica. 3.1.2 Estatísticas Básicas no Componente de Formação Geral A Tabela 3.5 apresenta as Estatísticas Básicas em relação ao componente da prova que avalia a Formação Geral dos estudantes concluintes.
Os alunos de todo Brasil obtiveram desempenho médio de 49,3.
Quanto à variabilidade, o Desvio padrão das notas dos estudantes do Brasil, como um todo, foi 16,5.
A maior Média foi obtida na região Sul ( 51,8 ), e a menor, na região Norte ( 41,5 ).
As demais médias foram : 47,1 na região Nordeste ; 50,1 na região Sudeste ; e 51,5 na região Centro-Oeste.
Já o maior Desvio padrão foi obtido na região Norte ( 17,3 ), e o menor, na região Centro-Oeste ( 15,3 ).
Os demais desvios padrões foram : 16,4 nas regiões Nordeste e Sudeste ; e 16,3 da região Sul.
A maior nota no Componente de Formação Geral da prova do Enade/2015 foi 95,2, obtida por, pelo menos, um aluno na região Centro-Oeste.
A menor nota Máxima foi obtida na região Norte ( 87,1 ).
A Mediana do Brasil, como um todo, foi 49,7, sendo a menor Mediana encontrada na região Norte ( 43,3 ) ; e a maior encontrada na região Centro-Oeste ( 53,0 ).
A nota Mínima nesta parte foi zero em todas as regiões, exceto na região Centro-Oeste ( 7,5 ).
Considerando -se as notas segundo Grande Região, observa -se que existe diferença estatisticamente significativa entre a menor Média do Componente de Formação Geral, obtida na região Norte ( 41,5 ), e as médias de todas as demais regiões.
O Gráfico 3.2 propicia a avaliação do desempenho dos estudantes do Componente de Formação Geral, a partir do histograma da distribuição das notas correspondentes.
A distribuição é unimodal, com moda nos intervalos ( 40 ; 60 ], enquanto na prova, como um todo ( Gráfico 3.1 ), a moda foi alcançada apenas no intervalo ( 50 ; 60 ].
Nota -se, ainda que, no Gráfico 3.2, as notas apresentam maior dispersão do que no Gráfico 3.1 ( distribuição das notas da prova ), confirmado pela comparação dos desvios padrões : 13,6 para a nota da prova como um todo e 16,5 para o Componente de Formação Geral.
Na Tabela 3.6 são apresentadas as informações referentes ao desempenho dos concluintes do Componente de Formação Geral, em diferentes agregações : Categoria Administrativa e Organização Acadêmica.
Observa -se que existe diferença estatisticamente significativa entre as médias dos tipos de Categoria Administrativa.
A maior Média foi obtida por estudantes de IES Públicas ( 54,6 ), e a menor, por estudantes de IES Privadas ( 49,0 ).
Considerando -se o tipo de Organização Acadêmica, nota -se que há diferença estatisticamente significativa entre suas médias. 3.1.3 Estatísticas Básicas do Componente de Conhecimento Específico A Tabela 3.7 apresenta as Estatísticas Básicas referentes ao Componente de Conhecimento Específico da área de Tecnologia em Gastronomia.
A Média do desempenho dos alunos do Brasil, como um todo, foi 52,2.
A maior Média foi obtida na região Sudeste ( 54,2 ), e a menor, na região Norte ( 44,2 ).
As demais médias foram : 49,4 na região Nordeste ; 53,6 na região Sul ; e 52,2 na região Centro-Oeste.
Quanto à variabilidade das notas, o Desvio padrão do Brasil, como um todo, foi 14,8, sendo o maior Desvio padrão observado na região Centro-Oeste ( 14,8 ), e o menor, na região Sul ( 14,1 ).
Os demais desvios foram : 14,7 nas regiões Norte e Sudeste ; e 14,5 na região Nordeste.
A Mediana das notas dos estudantes de todo o Brasil foi 52,4.
A maior Mediana ocorreu na região Sudeste ( 55,0 ), e a menor, na região Norte ( 44,4 ).
As demais medianas foram : 49,1 na região Nordeste ; 54,6 na região Sul ; e 51,6 na região Centro-Oeste.
A nota Máxima do Brasil, como um todo, foi 94,1, sendo obtida por, pelo menos, um aluno na região Nordeste.
As demais notas máximas foram : 79,2 na região Norte ; 92,6 na região Sudeste ; 89,0 na região Sul ; e 85,2 na região Centro-Oeste.
A nota Mínima foi zero em quase todas as regiões, exceto nas regiões Nordeste ( 3,7 ) e Centro-Oeste ( 7,4 ).
Observa -se que existe diferença estatisticamente significativa da menor Média das notas do Componente de Conhecimento Específico da região Norte ( 44,2 ) em relação às demais regiões.
O intervalo de confiança da região Norte ( 4,1 ) é significativamente maior do que de todas as outras.
Assim como os Gráficos 3.1 e 3.2, o Gráfico 3.3, apresentado a seguir, permite uma avaliação do desempenho de concluintes de Tecnologia em Gastronomia em relação ao Componente de Conhecimento Específico com um histograma da distribuição das notas correspondentes Esta também é uma distribuição unimodal, e o grupo modal é o ( 50 ; 60 ], o mesmo grupo modal da prova como um todo e para a Formação Geral ( junto com o intervalo ( 40 ; 50 ] ).
Assim como os Gráficos 3.1 e 3.2, o Gráfico 3.3, apresentado a seguir, permite uma avaliação do desempenho de concluintes de Tecnologia em Gastronomia em relação ao Componente de Conhecimento Específico com um histograma da distribuição das notas correspondentes Esta também é uma distribuição unimodal, e o grupo modal é o ( 50 ; 60 ], o mesmo grupo modal da prova como um todo e para a Formação Geral ( junto com o intervalo ( 40 ; 50 ] ).
A Tabela 3.8 apresenta uma comparação dos resultados em relação à Categoria Administrativa e à Organização Acadêmica, agora levando em conta o desempenho de alunos do Componente de Conhecimento Específico da prova.
No que se refere à Organização Acadêmica, a maior Média foi das Universidades ( 54,1 ), vindo a seguir a dos Centros Universitários ( 52,9 ) e, depois, a das Faculdades ( 49,5 ).
O maior Desvio padrão, e acima do valor para o Brasil, como um todo, foi o dos Centros Universitários ( 15,0 ).
As Universidades obtiveram a maior Mediana ( 54,5 ) e as Faculdades, a maior nota Máxima ( 94,1 ).
Por sua vez, os Centros Universitários obtiveram, respectivamente, 90,6 e 53,8 como nota Máxima e Mediana.
Já as Universidades obtiveram 92,6 como nota Máxima, enquanto as Faculdades obtiveram 49,4 como Mediana.
A nota Mínima foi zero para todas os Centros Universitários e as Faculdades, enquanto para as Universidades foi 3,7.
Quanto à Categoria Administrativa, observa -se um comportamento semelhante àquele da parte de Formação Geral e à prova como um todo, ou seja, existe diferença estatisticamente significativa entre as médias das IES Públicas ( 56,0 ) e IES Privadas ( 52,0 ).
Observa -se que existem diferenças estatisticamente significativas ao nível de 95 \% no Componente de Conhecimento Específico entre as notas das Faculdades em relação às Universidades e aos Centros Universitários. 3.2 ANÁLISE DAS QUESTÕES OBJETIVAS Esta seção apresenta estatísticas selecionadas e histogramas das Questões Objetivas de Formação Geral ( 3.2.1 ).
São também apresentadas e comparadas as médias das subpopulações caracterizadas por Grande Região. 3.2.1 Componente de Formação Geral A Tabela 3.9 apresenta as Estatísticas Básicas relativas às oito questões objetivas do componente da prova que abrange a Formação Geral dos estudantes.
A Média do Brasil foi 47,6.
A menor Média foi encontrada na região Norte ( 40,7 ), e a maior, na região Sul ( 51,5 ).
As demais médias foram : 44,9 na região Nordeste ; 48,3 na região Sudeste ; e 49,1 na região Centro-Oeste.
O Desvio padrão do Brasil foi 20,1, sendo o maior Desvio padrão encontrado na região Sul ( 20,3 ), e o menor, na região Centro-Oeste ( 18,5 ).
Os demais desvios foram : 19,4 na região Norte ; 20,2 na região Nordeste ; e 20,0 na região Sudeste.
As medianas ( 50,0 ), as notas máximas ( 100,0 ) e as notas mínimas ( 0,0 ) foram iguais para todas as regiões, exceto na região Norte, onde a Mediana e a nota Máxima foram, respectivamente, 37,5 e 87,5.
A Tabela 3.10 apresenta o Índice de Facilidade e o Índice de Discriminação ( ponto bisserial ) para cada uma das questões objetivas do Componente de Formação Geral.
Quanto ao Índice de Facilidade, foram usadas as seguintes cores para diferenciar o nível de dificuldade da questão :  Azul para as questões classificadas com índice Muito fácil ( \&gt ; =0,86 ), verde para as questões classificadas com índice Fácil ( 0,61 a 0,85 ), amarelo para as questões classificadas com Médio ( 0,41 a 0,60 ), vermelho para as questões classificadas com Difícil ( 0,16 a 0,40 ) e roxo para as questões classificadas com Muito difícil ( \&lt ; =0,15 ).
Já quanto ao Índice de Discriminação, foram usadas as seguintes cores para qualificar a questão :  As questões classificadas com índice Fraco receberam a cor vermelho ( \&lt ; =0,19 ), as classificadas com Médio receberam a cor amarelo ( 0,20 a 0,29 ), as classificadas com Bom receberam a cor verde ( 0,30 a 0,39 ) e as classificadas com Muito bom ( \&gt ; =0,40 ) receberam a cor azul.
As questões objetivas do Componente de Formação Geral, segundo o Índice de Facilidade, foram assim avaliadas : das oito questões, nenhuma teve o Índice de Facilidade classificado como Muito fácil.
Duas questões foram tidas como Fácil, por terem índice de acertos situado na faixa entre 0,61 e 0,85 ( de 61,0 \% a 85,0 \% de acertos ).
Três questões foram consideradas de dificuldade Médio, situando -se no intervalo entre 0,41 e 0,60 do Índice de Facilidade, ou seja, houve entre 41,0 \% e 60,0 \% de acertos, enquanto três questões foram classificadas na categoria Difícil, situando -se no intervalo entre 0,16 e 0,40.
Por fim, nenhuma das questões apresentou menos de 15 \% de acertos, razão pela qual seria classificada como Muito difícil.
Como já comentado, para análise das questões objetivas, relativas à Formação Geral, segundo o poder de discriminação, utilizou -se o Índice de Discriminação ( ponto bisserial ).
Nesta análise, as questões foram assim avaliadas : seis das oito questões apresentaram índice acima ou igual a 0,40 e, assim, foram classificadas com índice Muito bom para esse grupo de estudantes.
Duas questões tiveram Índice de Discriminação Bom, com valor entre 0,30 e 0,39, para esse grupo de estudantes.
Nenhuma questão teve nível Médio ou Fraco de discriminação para esse grupo de estudantes.
O Índice de Facilidade variou de 0,33 a 0,68, e o de Discriminação, de 0,31 a 0,46.
As seis questões com Índice de Discriminação Muito bom figuraram entre os diversos níveis de dificuldade desse conjunto : duas classificadas na categoria Fácil ( questões 5 e 6 ) do Índice de Facilidade, duas, na categoria Médio ( questões 3 e 8 ) e duas, na categoria Difícil ( questões 2 e 7 ).
Em particular, a questão 3 foi uma das que apresentou o maior poder discriminatório, com índice 0,46, porém foi considerada média em termos de facilidade, com uma proporção de 0,48 acertos.
O máximo de acertos foi alcançado pela questão 5 com um Índice de Facilidade de 0,68.
A questão de número 8 apresentou um Índice de Facilidade de 0,53, ou seja, pouco mais da metade dos estudantes conseguiu resolvê -la, dentro do universo de participantes.
Seu Índice de Discriminação foi Muito bom ( 0,41 ).
Já a questão 7 obteve Índice de Discriminação Muito bom, 0,46, porém seu Índice de Facilidade foi Difícil ( 0,34 ).
O Gráfico 3.4, para exemplificar, analisa o comportamento da questão de número 2 de Formação Geral.
Trata -se de uma questão considerada mediana, com relação à facilidade e a que obteve o maior Índice de Discriminação dessa parte da prova.
Neste gráfico, cada uma das cinco curvas representa o percentual de respostas em determinada alternativa da questão, em função do número de acertos dos estudantes nessa parte da prova ( Formação Geral/Múltipla Escolha ), antes de possíveis eliminações pelo critério do ponto bisserial.
A curva em preto corresponde à alternativa C, a correta para esta questão.
Assim, observa -se que entre os estudantes com menor número de acertos, nessa parte do exame, a situação mais frequente foi a escolha da alternativa D ( em roxo ), incorreta.
Por exemplo, entre os estudantes que acertaram três questões, 32,2 \% escolheram a alternativa D, 19,3 \% a alternativa C ( correta ), 20,3 \% a alternativa A, 14,2 \% a alternativa B e 12,8 \% a alternativa E.
Entre os que acertaram três respostas entre as questões de múltipla escolha de formação geral, 0,8 \% deixou a questão em branco e 0,4 \% marcou mais de uma alternativa, invalidando a questão.
À medida em que o número de acertos aumenta, indicando desempenho melhor nessa parte da prova, aumenta concomitantemente a proporção de estudantes que selecionaram a alternativa correta C, atingindo 100 \% para os estudantes com 8 acertos.
Essa análise permite verificar como a questão discriminou os grupos de desempenho, justificando o alto índice obtido na questão.
Cumpre notar que não é possível inferir deste gráfico nem o índice de facilidade ( que seria uma média da proporção ponderada pela quantidade de alunos com cada uma das notas ), nem o índice de discriminação ponto bisserial, por razão equivalente.
No caso extremo no qual a grande concentração dos acertos dos alunos fosse abaixo de 4, o índice de facilidade seria obrigatoriamente abaixo de 40 \% ( neste exemplo ).
Caso a concentração fosse em 6 acertos ou mais, o índice seria obrigatoriamente acima de 60 \%.
Os gráficos relativos às demais questões de Formação Geral constam do Anexo I. 3.2.2 Componente de Conhecimento Específico A Tabela 3.11 apresenta as Estatísticas Básicas em relação às questões objetivas do Componente de Conhecimento Específico da prova de Tecnologia em Gastronomia por Grande Região.
A Média do Brasil deste componente foi de 53,2.
A menor Média foi observada na região Norte ( 46,6 ), e a maior, na região Sudeste ( 54,7 ).
O Desvio padrão de todo o Brasil foi 14,8, sendo o menor Desvio padrão encontrado na região Sul ( 14,2 ), e o maior, na região Sudeste ( 14,8 ).
A Mediana de todo o Brasil foi 52,2, a mesma encontrada nas regiões Nordeste, e Centro-Oeste.
A região Norte apresentou valor menor para a Mediana ( 47,8 ), e as outras, valor maior : 56,5 para as regiões Sudeste e Sul.
A nota Máxima da prova foi 95,7, obtida nas questões objetivas do Componente de Conhecimento Específico, por, pelo menos, um aluno na região Nordeste, e, nas demais regiões, a nota Máxima da prova foi : 82,6 nas regiões Norte e Centro-Oeste ; 91,3 nas regiões Sudeste e Sul.
Em quase todas as regiões a nota Mínima foi zero, exceto nas regiões Nordeste ( 4,3 ) e Centro-Oeste ( 8,7 ).
A Tabela 3.12 apresenta os Índices de Facilidade e Discriminação ( ponto bisserial ) das questões objetivas do Componente de Conhecimento Específico para os estudantes de Tecnologia em Gastronomia.
Para facilitar a diferenciação das questões, usaremos as mesmas cores da Tabela 3.10 para as diferentes classificações dos Índices de Facilidade e de Discriminação.
Dentre as questões objetivas da parte da prova relativa ao Componente de Conhecimento Específico, nenhuma questão foi anulada pela Comissão.
Desse modo, a classificação, quanto ao Índice de Facilidade, foi estabelecida com base em todas as 27 questões.
A partir dos índices obtidos, pode -se concluir que um pouco menos da metade das questões objetivas da prova foi considerado, pelo menos, Difícil : das 27 questões, 12 foram classificadas como Difícil ( a classificação modal ) ou como Muito difícil.
Dez questões foram classificadas como Fácil ou como Muito fácil, e outras cinco consideradas como Médio.
Já quanto aos índices de discriminação das questões objetivas do Componente de Conhecimento Específico da prova, tem -se como resultado a seguinte classificação : 11 das 27 questões foram consideradas como boas, enquanto três delas tiveram Índice de Discriminação Muito bom.
Assim, para a maioria das questões – 14 em 27 – os Índices de Discriminação foram Bom ou Muito bom.
Dentre as demais, nove delas foram classificadas como Médio e outras quatro como Fraco, sendo 13, por conseguinte, a quantidade de questões nos dois patamares mais baixos de discriminação.
Constata -se, assim, que a prova – no que se refere ao Componente de Conhecimento Específico – possuía boa capacidade de discriminar entre aqueles que dominam ou não o conteúdo.
Dentre as questões que alcançaram os maiores índices de discriminação, três delas, as de números 20, 25 e 35, foram classificadas com Índice de Discriminação Muito bom, situando -se no intervalo de 0,40 a 0,43 do índice, todas elas foram classificadas na categoria Fácil, quanto ao Índice de Facilidade.
A questão de número 22 foi uma das mais difíceis dentre as 27 questões específicas, com baixo Índice de Facilidade, apenas 11,0 \% de acertos.
Essa questão apresentou poder discriminatório igualmente baixo, 0,03, o que comprova ter sido esta questão uma das mais difíceis para os estudantes.
Destaca -se, também, a questão 31, também com Índice de Facilidade 0,11, o que, em termos percentuais, corresponde a 11,0 \% de estudantes que responderam acertadamente.
Já o seu Índice de Discriminação foi 0,00.
Tais questões foram, portanto, pelo critério ponto bisserial, consideradas inadequadas.
Por isso, as questões 22 e 31 foram eliminadas do cômputo da nota final.
Além destas duas, as demais questões com Índice de Discriminação Fraco, questões 12 e 15 também não foram utilizadas no cômputo final das notas, num total de quatro questões eliminadas.
A título de exemplo das análises do comportamento das questões objetivas, o Gráfico 3.5 analisa a questão 35 do Componente de Conhecimento Específico.
Esta questão foi considerada pelos estudantes avaliados como uma questão de classificação Fácil como nível de facilidade da prova, apresentando Índice de Facilidade 0,79, ou seja, 79,0 \% dos estudantes assinalaram acertadamente a opção A, correspondente ao gabarito.
Seu Índice de Discriminação foi igual a 0,43, classificado como Muito bom, também sendo esta questão a que apresentou o maior índice discriminatório.
Neste gráfico, cada uma das cinco curvas representa o percentual de respostas em determinada alternativa da questão 35, em função do número de acertos dos estudantes nessa parte da prova, antes de possíveis eliminações de questões pelo critério do ponto bisserial.
A alternativa correta A, representada no gráfico pela curva em azul, foi escolhida em maiores proporções pelos alunos com desempenho melhor nessa parte da prova.
Já as alternativas incorretas, também denominadas distratores, foram selecionadas, principalmente, por aqueles com notas mais baixas.
Neste caso também a soma não é sempre 100 \% por causa das questões não respondidas ou com mais de uma opção marcada.
Aqueles com nota zero, na sua quase totalidade deixaram esta questão em branco ou marcaram mais de uma alternativa, comportamento considerado inválido.
A proporção de alunos que selecionou a resposta correta A aumenta gradativamente, chegando a atingir 100 \% para 19 acertos ou mais, enquanto a proporção dos que escolheram alternativas incorretas decai, a partir de dois acertos, como função do número de acertos nessa parte da prova.
Os gráficos relativos às demais questões do Conhecimento Específico constam do Anexo I. 3.3 ANÁLISE DAS QUESTÕES DISCURSIVAS Esta seção apresenta estatísticas selecionadas e histogramas das Questões Discursivas de Formação Geral ( 3.3.1 ) e Conhecimento Específico ( 3.3.2 ).
São também apresentadas e comparadas as médias de subpopulações caracterizadas por Grande Região. 3.3.1 Componente de Formação Geral As análises dos resultados de desempenho dos estudantes de Tecnologia em Gastronomia nas duas questões discursivas relativas à Formação Geral encontram -se na Tabela 3.13 e no Gráfico 3.6.
Na Tabela 3.13, observa -se que a nota Média nesse conjunto de questões foi próxima da obtida nas objetivas.
Os estudantes, de todo o Brasil, obtiveram, em Formação Geral, Média 47,6 nas questões objetivas e 51,8 nas questões discursivas.
No entanto, pode -se notar um aumento do Desvio padrão de 20,1 nas questões objetivas do Componente de Formação Geral dos alunos de todo o Brasil, para 22,8 nas questões discursivas do mesmo componente.
A maior Média foi obtida na região Centro-Oeste ( 55,1 ), e a menor, na região Norte ( 42,7 ).
A Mediana de todo o Brasil e da região Sul, neste componente, foi 58,0, na região Sudeste a Mediana foi maior, 58,5.
Nas demais regiões foi menor : Norte ( 50,0 ), Nordeste ( 57,0 ) e Centro-Oeste ( 60,5 ).
A nota Máxima ( 95,0 ) foi obtida na região Nordeste, enquanto a menor nota Máxima foi obtida na região Centro-Oeste ( 88,0 ).
A nota Mínima ( 0,0 ) foram as mesmas em todas as regiões do Brasil, sem exceção.
O Gráfico 3.6 representa a distribuição das notas nas questões discursivas do Componente de Formação Geral.
A moda desta distribuição ocorre no intervalo ( 60 ; 70 ].
Destaca -se, também, o intervalo [ 0 ; 10 ], um máximo local com distribuição próxima a 10 \% do total de notas, sendo que no intervalo [ 0 ; 10 ] se inclui, além da nota zero, a frequência de alunos que deixaram este tipo de questão em branco.
Na sequência, os resultados verificados para cada uma das questões discursivas de Formação Geral serão apresentados, estabelecendo -se relações com os conteúdos abordados em cada uma delas.
Os comentários da Banca de docentes corretores, a respeito do observado na correção das respostas dos estudantes, suas impressões e conclusões serão apresentados junto à análise de cada questão.
Cumpre esclarecer que, tendo em vista que as questões discursivas de Formação Geral são padronizadas, ou seja, constam de todas as provas, os comentários da Banca são os mesmos para todas as carreiras acadêmicas, sendo direcionados a todos os estudantes que participaram do Enade/2015.
A seguir, serão analisados os desempenhos da Área de Tecnologia em Gastronomia nas duas questões discursivas de Formação Geral do Enade/2015, comparando -se os resultados obtidos com comentários para cada questão. 3.3.1.1 Análise de Conteúdo da Questão Discursiva 1 do Componente de Formação Geral Os dados de Tecnologia em Gastronomia, obtidos a partir das respostas à questão 1, encontram -se na Tabela 3.14 e no Gráfico 3.7.
Nessa questão – de melhor desempenho dentre as duas de Formação Geral – os alunos, de todo o Brasil, tiveram Média 56,6.
A maior Média para a questão 1 foi obtida na região Centro-Oeste ( 58,4 ), e a menor, na região Norte ( 46,2 ).
Quanto à variabilidade das notas, o Desvio padrão de todo o Brasil foi 25,2.
O menor Desvio padrão foi obtido na região Sudeste ( 24,2 ), e o maior Desvio padrão foi obtido na região Norte ( 30,6 ).
As medianas das regiões Sudeste, Sul, Centro-Oeste e do Brasil, como um todo, foram iguais ( 65,0 ).
Na região Norte, a Mediana foi 55,0 ; e na região Nordeste, foi 60,0.
As notas máximas ( 100,0 ) foram obtidas nas regiões Sudeste e Sul.
Na região Norte, a nota Máxima foi 90,0, e nas regiões Nordeste e Centro-Oeste, foi 95,0.
A nota Mínima ( 0,0 ) foi a mesma para todas as regiões do Brasil, sem exceção.
O Gráfico 3.7 mostra a distribuição das notas na questão discursiva 1 do Componente de Formação Geral.
Observa -se que a maior frequência corresponde aos alunos que obtiveram nota dentro do intervalo ( 70 ; 80 ], correspondendo à moda da distribuição.
Destacam -se, também, os alunos que deixaram a questão em branco, representando um máximo local com mais de 10 \% do total. 3.3.1.2 Comentários sobre a correção de Conteúdo das respostas à Questão Discursiva 1 O enunciado da Discursiva 1 de Formação Geral estava claro e indicava um tema de extrema relevância, qual seja : a busca pela educação universal, especificamente, que apontava um caso onde ocorreu a discriminação de acesso ao ensino formal atingindo uma menina, representando a categoria das mulheres discriminadas em tantos campos da vida e, mais ainda, em sociedades que cultuam o poder soberano dos homens.
A linguagem utilizada foi clara, indicando a existência dos fatos vinculados à luta da menina Malala pelo acesso ao ensino formal das jovens do seu país, e o confronto ideológico que gerou o ataque à sua integridade física.
Diante da repercussão internacional, houve o reconhecimento pelo prêmio Nobel, e a ONU resolveu apoiá -la como símbolo da luta pela igualdade entre os sexos buscando a universalização da educação.
O conteúdo está integrado às Diretrizes, vez que o acesso ao ensino formal está vinculado ao direito constitucional à educação, que, também, é protegido pelo Estatuto da Criança e do Adolescente e nos pactos de Direitos Humanos aderidos pelo Brasil.
Por outro lado, o texto demandava uma reflexão sobre a realidade dos indivíduos que não têm acesso à educação e a situação das mulheres excluídas desse processo por segmentos da sociedade.
Como se tratava de texto contendo uma ampla gama de possibilidades, era natural a ocorrência de interpretações diversas das esperadas.
Releva notar que o enunciado pretendeu, primordialmente, identificar a questão de gênero, apontando para as dificuldades que as mulheres possuem, não somente no Brasil, mas em vários países do mundo.
Ao analisar as provas realizadas, constatamos que no item ‘ a ’ poucos foram os que identificaram os estatutos normativos indicados no padrão de resposta.
Essa ausência pode ser justificada tendo em vista que, no enunciado \textbf{indicado}, não existia elemento que pudesse ser utilizado pelo concluinte para embasar sua resposta nesse caminho.
As respostas foram mais gerais, identificando que a educação constitui um item fundamental para o desenvolvimento do indivíduo e que as mulheres sofrem mais com as limitações impostas por países com regime patriarcal, que discriminam o sexo feminino e, com frequência, proíbem as mulheres de frequentar a escola.
Deve -se ter em vista que a indicação de resposta que vinculava o texto aos estatutos normativos representou um raciocínio próprio dos estudantes de Direito que, naturalmente, encaminharam suas respostas para a fundamentação nessa linha identificando, não somente a Constituição, mas também leis, regulamentos, etc.
Tal linha de argumentação não foi objeto da formação dos estudantes das demais áreas.
A maior parte dos concluintes estabeleceu o tema educação formal associado à busca pelo conhecimento, como acesso ao ensino, indicando a independência intelectual e o desenvolvimento da consciência crítica, limitando as manipulações e permitindo amplo acesso ao mercado de trabalho, traduzindo uma libertação econômica e ascensão social.
No tocante ao item ‘ b ’, restou claro que a maioria esmagadora das respostas confluiu para as indicações do padrão apresentado, sendo relevante o tema de igualdade de gênero, além de candente para as concluintes do sexo feminino que desenvolveram os textos de forma concatenada, fundamentada e, muitas vezes, lamentando que, nos dias atuais, ainda se discuta a necessidade de igualar os gêneros na sociedade.
Nessa linha, foram identificadas respostas, consoante com o padrão apresentado, especialmente nos seguintes temas :  violência física e psicológica contra a mulher, incluindo a Lei Maria da Penha, no caso específico do Brasil ;  tolerância/intolerância a vestimentas, trajes, comportamentos socialmente estereotipados ;  aspectos socioculturais que impõem à mulher uma condição de submissão na sociedade, tais como a impossibilidade de a mulher manifestar seus desejos e posicionamentos em algumas culturas, entre outros ;  igualdade/desigualdade de gênero, por exemplo no mercado de trabalho, em relação à desigualdade salarial ;  ideais de Liberdade, Igualdade e Fraternidade : referência a esses ideais como possibilidade de equilibrar as relações de poder entre homens e mulheres.
Em muitos casos, também, os concluintes fizeram um liame entre o acesso à educação e a melhora das relações entre homens e mulheres ao permitir que a conclusão de maior grau de educação formal permita o acesso a cargos públicos relevantes, empregos bem remunerados e posições de poder, tanto na área pública quanto na área privada.
Como as opções do segundo item foram mais amplas, o padrão de resposta identificou as linhas de pensamento adotadas pelos concluintes.
A maior parte das respostas identificou, claramente, o problema da relação de gênero e suas consequências na sociedade.
Também foi frequente constatarem que existe efetiva melhora na situação das mulheres no Brasil, mas que tal movimento não ocorre em todo o mundo.
Aliás, o exemplo de Malala conduziu a uma reflexão em âmbito mundial, e os concluintes apresentaram suas análises críticas à situação da mulher em países que não admitem sequer a sua inserção no plano educacional e, muito menos, no mercado de trabalho.
A maior parte dos concluintes apresentou concatenação lógica nas suas respostas e uma minoria, pouco expressiva, não demonstrou esse domínio de desenvolvimento lógico dos textos.
Ressalte -se que essa referência não inclui a utilização do padrão culto da língua portuguesa que tem avaliação própria. 3.3.1.3 Análise de Conteúdo da Questão Discursiva 2 do Componente de Formação Geral A Tabela 3.15 mostra que o desempenho médio dos estudantes na questão discursiva 2 ( média 49,7 ) foi inferior ao obtido na questão discursiva 1 ( média 56,6 ).
A região Centro-Oeste foi aquela cuja Média, nessa questão, foi maior ( 53,0 ), e a de menor Média foi a região Norte ( 40,5 ).
Quanto à variabilidade das notas, o Desvio padrão de todo o Brasil foi 25,5, superior ao obtido na questão discursiva 1 ( 25,2 ).
O maior desvio nessa questão foi obtido na região Norte ( 29,7 ), enquanto o menor foi obtido na região Centro-Oeste ( 24,5 ).
A Mediana de todo o Brasil foi 59,0, a mesma das regiões Nordeste, Sudeste e Sul, enquanto a maior Mediana foi encontrada na região Centro-Oeste ( 60,0 ) e a menor, na região Norte ( 47,0 ).
As notas máximas ( 98,0 ) foram as mesmas nas regiões Nordeste, Sudeste e Sul, enquanto nas demais regiões foram 95,0.
As notas mínimas ( 0,0 ) foram as mesmas em todas as regiões do país, sem exceção.
O Gráfico 3.8 mostra a distribuição das notas na questão discursiva 2 do Componente de Formação Geral.
Observa -se que a maior frequência corresponde aos alunos que obtiveram nota situada no intervalo ( 60 ; 70 ].
Nota -se, ainda, que as notas ficaram mais distribuídas ao longo dos intervalos em comparação à questão discursiva de número 1. 3.3.1.4 Comentários sobre a correção de Conteúdo das respostas à Questão Discursiva 2 O enunciado da Discursiva 2 estava, igualmente, claro, indicando um tema de extrema relevância, qual seja : o reconhecimento de manifestações artísticas nascidas na periferia, como integrante da cultura reconhecida pela sociedade.
A linguagem utilizada foi clara, indicando a análise dos fatos e suas consequências no cotidiano dos partícipes e o confronto com outras visões culturais que são consideradas mais adequadas pela sociedade civil.
O conteúdo está integrado às Diretrizes, vez que o aspecto cultural e sua inter-relação com o conhecimento haurido no nível superior é essencial para a compreensão de mundo do universitário concluinte, com reflexo na sua atuação profissional.
O texto demanda uma reflexão sobre a realidade dos indivíduos que não têm acesso aos itens básicos da sociedade, como educação, saúde, cultura tradicional e segurança e que, apesar disso, criam alternativas para o seu lazer.
A incorporação do funk como manifestação cultural e artística da sociedade brasileira encontrou, basicamente, duas correntes de resposta :  os que aquiescem positivamente à integração ;  os que rejeitam, acerbamente, tal possibilidade.
No meio, uma corrente que somente admite a integração com modificações na conduta dos integrantes do movimento, com a exclusão de letras exclusivamente com apelos sexuais e incitação à violência.
A interpretação positiva indica a origem da música na periferia, notadamente nas comunidades cariocas e sua projeção para o Brasil e o mundo.
Nessa perspectiva, surge o reconhecimento como cultura popular comparada ao samba, axé e outros ritmos também nascidos no Brasil.
Também ocorre comparação com o movimento hip-hop americano.
Já a análise negativa prioriza a vinculação do ritmo com traficantes, prostitutas, vilipêndio às mulheres, falta de controle quanto ao som, local dos eventos, transações sexuais em público, uso frequente de drogas, violência exacerbada.
Faltam, para essa corrente, elementos básicos da cultura, tratando -se de música sem autonomia como movimento de mudança ou aspecto próprio.
A corrente intermediária reconhece o vínculo originário com as comunidades dominadas por quadrilhas de traficantes, mas admite o seu reconhecimento de acordo com a mudança dos atores que transitam na modalidade, como músicos, intérpretes, para adequar as letras das canções expurgando os aspectos considerados nocivos como o apelo sexual, as indicações positivas para a violência e apologia às drogas.
Um aspecto bastante criticado nas músicas dessa modalidade foi a exploração da mulher, entendida como um objeto de satisfação masculina.
É relevante tal constatação, pois a questão pertinente à Malala reflete a luta pela valorização da mulher.
Os concluintes compreenderam a ideia que foi problematizada e escolheram uma estratégia adequada para solucioná -la.
O caráter dinâmico da cultura foi comumente expressado, não através do enfrentamento abordado no texto, mas mediante a indicação ( de forma rasa e genérica ) das dificuldades e consequente aceitação social de manifestações culturais, antes marginalizadas e que hoje integram o acervo cultural nacional, tais como : samba, grafite, capoeira e forró.
Quanto às reflexões de Laraia, a grande maioria reconheceu o funk como manifestação cultural, porém, o identificou mais na sua expressão musical, no balanço de seu ritmo, tecendo contundentes críticas, e mesmo forte discordância quanto ao conteúdo e mensagem das letras quando remetem à depreciação da imagem feminina ou induzem à apologia ao crime.
De resto, os textos abordaram bem mais o item dois do padrão de Respostas.
Uma minoria enfrentou os temas mudança e dinamismo cultural.
As respostas da corrente positiva indicam que o preconceito sempre acompanhou as manifestações populares como o samba, na sua origem.
Admitem que diante da miscigenação surgem movimentos culturais de várias naturezas, o que inclui o funk.
Também indicam a dificuldade de lazer dos membros das comunidades o que faz surgirem eventos para preencher essa necessidade.
Nesse âmbito, apontam a música como alternativa de ascensão social e integração aos bens de consumo.
As respostas da corrente negativa vinculam a música ao tráfico de drogas cujos elementos comandam as comunidades, propiciando que os bailes sejam mecanismos de venda e consumo dessa mercadoria ilícita.
Na mesma toada, indicam a vinculação com a exploração sexual, principalmente das mulheres maiores e menores de idade.
Indicam aspectos imorais nas letras.
Apresentam os locais dos bailes como inadequados e provocadores da desordem pública.
Recusam que músicas com discriminação às mulheres sejam consideradas como cultura.
A corrente intermediária indica a possibilidade de adoção do funk como cultura desde que ocorram modificações nas letras e no comportamento dos atores ( músicos, letristas e público ), retirando os aspectos negativos já apontados.
Fala -se em “ funk do bem ”.
A maior parte das respostas apresentadas indica que os concluintes têm conhecimento dos conflitos sociais e, no caso da questão do funk, conseguem identificar a origem popular da música.
A partir daí é que ocorre a cisão entre os que defendem o movimento como cultura, aduzindo a ascensão social dos artistas envolvidos dos que vinculam a música com atos ilícitos, rejeitando sua inserção no âmbito cultural.
Em algumas provas, há dificuldade de concatenação lógica de raciocínio, mas a maior parte consegue desenvolver o texto de forma razoável.
Essa questão permite \textbf{visualizar} a divisão social com os pensamentos liberais e conservadores, sendo bastante claros nas respostas dos concluintes. 3.3.1.5 Análise de Língua Portuguesa das Questões Discursivas do Componente de Formação Geral Os dados de Tecnologia em Gastronomia, obtidos a partir das respostas às questões discursivas do Componente de Formação Geral, no que tange à Língua Portuguesa, encontram -se na Tabela 3.16 e no Gráfico 3.9.
Nesse aspecto, os alunos, de todo o Brasil, tiveram Média 51,2.
A maior Média com respeito à Língua Portuguesa foi obtida na região Centro-Oeste ( 54,7 ), e a menor, na região Norte ( 42,3 ).
Quanto à variabilidade das notas, o Desvio padrão de todo o Brasil foi 23,2.
O menor Desvio padrão foi obtido na região Centro-Oeste ( 22,2 ) e o maior Desvio padrão foi obtido na região Norte ( 27,2 ).
A Mediana das notas de Língua Portuguesa foi 57,0 para o Brasil como um todo, sendo a maior Mediana e a menor Mediana obtidas, respectivamente, nas regiões Centro-Oeste ( 60,0 ) e Norte ( 49,0 ).
A nota Máxima para todo o Brasil foi de 98,0, com, pelo menos, um aluno tirando essa nota na região Nordeste.
Além disso, a nota Mínima foi zero em todas as regiões do país, sem exceção.
O Gráfico 3.9 mostra a distribuição das notas de Língua Portuguesa do Componente de Formação Geral.
Observa -se que a maior frequência ( quase 25 \% ) corresponde aos alunos que obtiveram nota no intervalo ( 60 ; 70 ]. 3.3.1.6 Comentários sobre a correção das respostas de Formação Geral com respeito à Língua Portuguesa Ao encaminhar as questões 1 e 2 na direção da produção de um texto dissertativo, espera -se que o participante utilize seus conhecimentos sobre o assunto e estruture seus textos de acordo com as exigências do registro formal próprio dessa situação comunicativa.
Essa configuração determina exigências quanto : à adequação da seleção vocabular, ao desenvolvimento do conteúdo, à estruturação sintática dos períodos, à organização lógica das ideias, à utilização de procedimentos de encadeamento textual e referenciação, à obediência às exigências morfossintáticas próprias da modalidade escrita da norma-padrão, ao respeito às regras ortográficas e às regras de acentuação gráfica.
O padrão de resposta utilizado na avaliação das questões 1 e 2 considerou os aspectos relevantes ao bom desempenho linguístico como competências distintas, de modo a permitir um mapeamento detalhado do domínio dos recursos disponíveis na Língua Portuguesa para a comunicação escrita formal.
Com base nesse objetivo, foram avaliados os seguintes aspectos : a ) Estruturação textual condizente com o gênero solicitado ( texto dissertativo ) e o modo de organização textual expositivo adequado ao gênero – essa competência envolve :  distribuição do conteúdo do texto em parágrafos, de modo a garantir a sua organização temática ;  estruturação sintática condizente com o padrão da modalidade escrita formal da língua portuguesa de modo a garantir a clareza necessária ;  utilização de operadores discursivos que contribuam para a progressão temática do texto, estabelecendo relações lógicas entre as ideias apresentadas, tanto do ponto de vista intrafrasal, como do interfrasal ;  utilização de procedimentos de referenciação lexical e pronominal que permitam a retomada de referentes textuais ;  utilização de sinais de pontuação que contribuam para a organização lógica da frase e do texto.
Espera -se, portanto, que o participante recorra a procedimentos linguístico-discursivos para organizar seu texto, permitindo o encadeamento lógico entre suas partes de forma a garantir a progressão e a coerência textual.
Isso significa que os seguintes procedimentos foram penalizados, de acordo com o padrão de resposta proposto :  elaboração de frases fragmentadas que comprometam a estrutura lógico-gramatical do texto ;  sequência justaposta de ideias sem encaixamentos sintáticos, reproduzindo hábitos da oralidade ;  elaboração de frase com apenas oração subordinada, sem oração principal ;  emprego equivocado do conector ( preposição, conjunção, pronome relativo, alguns advérbios e locuções adverbiais ) que não expresse a relação lógica adequada entre dois trechos do texto e prejudique a compreensão da mensagem ;  repetição ou substituição inadequada de palavras sem se valer dos recursos oferecidos pela língua ( pronome, advérbio, artigo, sinônimo ) ;  utilização inadequada dos sinais de pontuação comprometendo a clareza textual. b ) Respeito às convenções ortográficas da norma-padrão da Língua Portuguesa – essa competência envolve o domínio das regras de acentuação gráfica e da grafia padrão das palavras ( com ausência de abreviaturas próprias da linguagem da internet ), de acordo com as convenções estabelecidas pela legislação em vigor e consubstanciadas no Vocabulário Ortográfico da Língua Portuguesa, editado pela Academia Brasileira de Letras ( com aceitação da legislação anterior, no caso das regras relativas ao uso do hífen e da acentuação gráfica, já que o exame foi realizado no mês de novembro de 2015 ).
A competência relativa ao respeito às convenções ortográficas envolve :  grafia correta das palavras ;  respeito às regras de acentuação gráfica ;  emprego de maiúsculas em início de frase, em nomes próprios de pessoas, lugares ou instituições ;  ausência de abreviações como p/, vc, tb, pq, tá, né, usadas muitas vezes em escrita informal e na internet ;  obediência às regras de separação de sílabas no final da linha. c ) Domínio dos diferentes aspectos morfossintáticos próprios da modalidade escrita formal da norma-padrão da Língua Portuguesa – os itens relacionados a essa competência são : a concordância nominal, a concordância verbal, a regência nominal, a regência verbal, a flexão nominal, a flexão verbal, a correlação entre os tempos verbais, a colocação pronominal.
O domínio dessa competência envolve :  flexão do verbo para estabelecer concordância de número com o sujeito da frase ;  flexão do artigo, do adjetivo e do pronome para concordar em número e em gênero com o substantivo a que se referem ;  respeito à regência nominal e à verbal, utilizando a preposição adequada depois de um substantivo, de um verbo ou de um adjetivo ; emprego da preposição antecedendo o pronome relativo e emprego adequadamente os pronomes relativos “ cujo(a) ” e “ onde ” ;  emprego adequado do acento grave indicador da crase entre uma preposição e um artigo ( a+a ) ;  atendimento às regras de colocação pronominal ( próclise e ênclise ), distintas dos hábitos da oralidade ou da escrita informal. d ) Seleção vocabular adequada à modalidade escrita formal da Língua Portuguesa, exigida pela situação comunicativa – essa competência envolve a precisão na utilização do vocabulário relacionado à temática solicitada pela questão ; a ausência de marcas da oralidade, como termos de sentido muito genérico ( “ coisa ”, “ negócio ”, “ você ” ) e termos de registros mais informais ( como gírias, jargões, frases feitas, ditados populares, termos regionais ).
Assim, espera -se que o participante respeite a adequação vocabular não empregando gírias ou expressões coloquiais, evite repetição desnecessária de palavras e utilize um vocabulário mais formal, como solicitado por um texto dissertativo.
A escolha dessas competências para subsidiar o processo de avaliação apoia -se na concepção de que, no desempenho dos graduandos, a modalidade escrita tem apresentado uma intensa simplificação, originada no padrão da modalidade oral da Língua Portuguesa.
No caso do texto de base dissertativa, inscrito em um registro formal, a distância entre as duas modalidades é ainda maior, o que provoca situações de hipercorreção ( desvios provocados pela utilização inadequada de uma regra da norma-padrão ) e de truncamentos sintáticos ( estruturas frasais incompreensíveis devido à complexidade sintática própria da modalidade escrita ).
Observam -se, então, os seguintes aspectos que marcam essa distinção entre as duas modalidades, devido à excessiva simplificação da modalidade falada : a ) redução drástica de estruturas subordinadas, compensada pelo aumento na frequência de estruturas coordenadas e absolutas, por um lado, ou pela elaboração de estruturas truncadas pelo excesso de ideias sem a devida conexão subordinativa ; b ) redução no uso de operadores argumentativas para expressar relações lógicas essenciais à construção do texto, substituídas pela exigência de inferência por parte do interlocutor para suprir a sua ausência ; c ) redução do uso do subjuntivo, ao lado da ampliação do uso do indicativo combinado a estruturas frasais coordenadas ou absolutas ; d ) empobrecimento do processo de referenciação, com a repetição de pronomes ou nomes ; e ) simplificação extrema da marcação da categoria tempo na morfologia verbal ; f ) falta de domínio de vocabulário mais abstrato e de maior complexidade, essencial ao desenvolvimento do processo dissertativo ; g ) redução drástica no emprego da acentuação gráfica, processo intensificado pela intensa utilização das redes sociais.
Os aspectos macroestruturais da elaboração do texto não foram avaliados neste processo porque dizem respeito à avaliação do conteúdo.
Portanto, a banca de Formação Geral, composta por profissionais de diferentes áreas do conhecimento, se encarregou da avaliação do atendimento ao solicitado no enunciado das questões do ponto de vista do desenvolvimento do conteúdo.
A avaliação do desempenho linguístico considerou, portanto, três grandes grupos de competências, segundo os aspectos explicitados anteriormente : a ) domínio das convenções ortográficas : grafia de vogais e consoantes, uso de maiúsculas e minúsculas, emprego do hífen e acentuação gráfica ; b ) domínio dos procedimentos de estruturação textual do ponto de vista microestrutural : organização interna dos períodos, emprego de conectores para a articulação lógica entre os períodos e entre os parágrafos, emprego de marcas de referenciação lexical e pronominal ; utilização dos sinais de pontuação que contribuem para a organização lógica da frase. c ) domínio das regras de caráter morfossintático estabelecidas como modelares do ponto de vista da modalidade escrita formal da norma-padrão da Língua Portuguesa : concordância nominal e verbal, regência nominal e verbal, colocação pronominal, flexão nominal e verbal, correlação entre tempos e modos verbais, ausência de marcas de oralidade.
A seleção vocabular adequada à modalidade escrita formal da Língua Portuguesa foi incorporada à essa última competência, tendo em vista a intersecção entre as duas do ponto de vista das exigências do registro formal da modalidade escrita da norma-padrão.
Apresenta -se, a seguir, o padrão de resposta aprovado pelo Inep e já utilizado na avaliação do desempenho linguístico das questões 1 e 2 de Formação Geral no Enade/2013 e no Enade/2014.
Às competências, reunidas nos três grupos descritos anteriormente, foram atribuídos os seguintes pesos relativos : aspectos ortográficos ( 20 \% ) ; aspectos textuais ( 40 \% ), aspectos morfossintáticos e vocabulares ( 40 \% ).
Com base na avaliação do processo desenvolvido nos dois anos anteriores, utilizaram -se quatro níveis para a avaliação do desempenho linguístico nas três competências.
O nível zero ficou reservado para casos especiais, como fuga total ao tema, palavras soltas, produção de uma frase incompleta, entre outros.
A avaliação das respostas elaboradas pelos participantes revelou um resultado coerente na comparação entre as duas questões.
Entretanto, em relação aos resultados do Enade/2014, observou -se uma melhora do domínio de todas as competências, talvez devido ao perfil dos cursos analisados, como Direito e Jornalismo, que pressupõem a utilização profissional da linguagem.
Observou -se uma diferença no desempenho dos participantes em relação às duas questões discursivas, provavelmente em função dos temas solicitados : - a Questão 1 solicitava que o participante desenvolvesse um texto dissertativo sobre o significado da premiação de Malala, abordando dois aspectos : o direito das jovens à educação formal e as relações de poder entre homens e mulheres no mundo.
Esse encaminhamento favoreceu o desenvolvimento mais articulado do texto, o que se refletiu em melhor desempenho dos aspectos vocabulares e textuais. - a Questão 2 solicitava que o participante desenvolvesse um texto dissertativo para discutir a questão do reconhecimento do funk como legítima manifestação artística e cultural do povo brasileiro.
Por ser um tema muito polêmico, o tema do preconceito sócio-cultural-econômico favoreceu a construção de textos mais críticos, dividindo o conjunto de participantes em dois grupos : os que defendiam esse tipo de manifestação artística como legítima e os que a criticavam de maneira veemente, chegando, inclusive, a utilizar muitos termos de baixo calão.
Essa última atitude teve como consequência a elaboração de textos mal desenvolvidos do ponto de vista sintático-discursivo, com muitos truncamentos.
Os resultados revelaram essa diferença.
Ao analisar os dados de cada questão, constata -se que a Questão 1 teve maior número de questões de médio ou bom desempenho, em relação à Questão 2.
Por outro lado, a Questão 2 apresentou um índice maior de notas baixas em relação à Questão 1.
Os resultados da avaliação correspondem aos seguintes aspectos observados em cada competência : a ) aspectos ortográficos : o desempenho dos participantes revelou uma diferença muito grande nos dois aspectos analisados : baixo índice de desvios ortográficos e grande índice de desvios de acentuação.
Em vários casos, ocorre ausência completa de acentuação gráfica.
Os resultados revelam que a tendência dominante entre os universitários brasileiros é a eliminação da acentuação gráfica, talvez motivada pelos hábitos relacionados às redes sociais e pela ausência de esclarecimento dos meios de comunicação, das autoridades e das escolas sobre as decisões do Acordo Ortográfico da Língua Portuguesa de 1990.
Os casos mais sistemáticos de eliminação do acento indicador da sílaba tônica são :  palavras proparoxítonas ( “ proximos ”, “ politicos ” ) ;  palavras paroxítonas terminadas em ditongo crescente ( “ necessario ”, “ noticiarios ”, “ individuo ”, “ dependencia ”, “ varios ”, “ propria ” ) ;  palavras oxítonas ( “ ninguem ”, “ esta ”, “ ate ”, “ tambem ”, “ alvara ” ) ;  uso indevido do acento gráfico em palavras como “ genêro ”, “ melâncolia ”, por exemplo.
Quanto ao domínio das convenções relativas à grafia das palavras, observam -se desvios como : I. hipercorreção pela escolha de “ e ” no lugar de “ i ” ou de “ i ” no lugar de “ e ”, por influência de hábitos da oralidade ( “ descriminação ” por “ discriminação ”, “ entervenção ” por “ intervenção ”, “ esteriotipada ” por “ estereotipada ” ) ; II. eliminação do “ r ” marcador do infinitivo verbal ( “ esta ” no lugar de “ estar ”, “ estuda ” no lugar de “ estudar ” ).
III. eliminação de uma sílaba na palavra “ educação ”, que foi grafada “ edução ” por um grande número de participantes.
IV. outros casos de desvios de grafia relacionados à variação diastrática podem ser observados em “ estrupo ”, “ subjulgadas ”, “ precoseito ”, “ soubre ”, “ vecendo ”, “ indesencia ”, “ apolojia ”, “ fulga ”.
V. desvios de segmentação podem ser atribuídos à confusão entre o morfema gramatical “ mos ” e o pronome oblíquo “ nos ”, em “ chegar mos ”, “ enfrentar -mos ”.
VI. inadequação no uso da maiúscula para destacar determinadas palavras-chave do texto, como “ Violência ”, “ Brasileiros ”.
Destaque -se, também, o grande número de participantes que grafam os textos inteiramente em caixa alta.
Vale observar, também, que, ao contrário do que se esperava, não apareceram abreviaturas próprias do “ internetês ”, relacionado ao uso de redes sociais e emails. b ) aspectos textuais : esta competência é a que se revela como a mais problemática entre os participantes, porque são muitos os problemas observados, desvios acumulados durante toda a formação do estudante : sequência justaposta de ideias sem encaixamentos sintáticos ; redução drástica de estruturas subordinadas, ao lado do aumento na frequência de estruturas coordenadas e absolutas ; redução no uso de conectores para expressar relações lógicas essenciais à construção do texto, substituídas pela exigência de inferência por parte do interlocutor para suprir a sua ausência ; emprego equivocado de operadores que não estabelecem relações lógicas coerentes entre ideias do texto ; emprego inadequado do pronome relativo ( com omissão da preposição ou a utilização de pronome inadequado, como “ onde ” ) ; repetição exaustiva de termos sem a utilização de procedimentos mais sofisticados de substituição ( hiperonímias, hiponímias, nominalizações, expressões metafóricas ) ; frases fragmentadas que comprometem a estrutura lógico-gramatical ; frases formadas apenas por oração subordinada, sem oração principal.
Um importante aspecto a destacar é o baixíssimo desempenho de uma parte dos participantes em relação à estrutura formal do texto produzido, o que é extremamente preocupante ao se levar em conta que são graduandos em fase final de formação.
São frequentes os casos de desvios de estruturação frasal, com uso inadequado ou ausência de conectivos entre parágrafos e entre frases.
Em uma parte dos textos, falta um mínimo de textualidade e de domínio do registro padrão da língua.
Na verdade, observam -se relações linguísticas quase agramaticais, como as estabelecidas pela sequência de gerúndios sem o apoio de um ponto de partida para a organização das informações gramaticais e semânticas.
Observou -se que uma grande parte dos participantes não organizou as ideias em parágrafos, talvez devido ao pequeno número de linhas disponibilizadas para a resposta da questão ou, quem sabe, pela suposição de que não seria necessária essa divisão por não se tratar de um texto no modelo de uma redação dissertativo-argumentativa, como solicitado nos vestibulares.
Quanto à utilização dos mecanismos de referenciação, deve -se destacar a ocorrência, em uma boa parte dos textos, de repetições de palavras ou expressões sem a utilização de termos sinônimos ou pronomes, como seria adequado.
Outro aspecto relevante a ser destacado é a quase total ausência de operadores argumentativos, tanto intrafrasais como interfrasais, repercutindo uma tendência atual da mídia escrita.
Quanto à utilização dos sinais de pontuação, observou -se uma grande precariedade nos textos analisados.
É muito frequente a ocorrência de parágrafos sem marca interna de pontuação para separar os períodos.
Vale observar que não foi penalizada a ausência de vírgula para destacar locuções ou adjuntos adverbiais de pequena extensão deslocados de posição na frase, por ser um uso opcional.
São os seguintes os tipos de problemas encontrados : I. vírgula : utilização de vírgula para separar o sujeito e o predicado ; ocorrência de apenas uma das vírgulas para separar uma palavra, uma expressão ou uma oração encaixada ; uso de vírgula no lugar do ponto para separar ideias que constituem períodos distintos ; II. ponto e vírgula : utilização do ponto e vírgula no lugar de vírgula ; III. ponto final : ausência de ponto final para separar períodos. c ) aspectos morfossintáticos e vocabulares : os resultados são muito transparentes em relação aos aspectos mais problemáticos no desempenho dos participantes.
O desvio mais frequente, em relação à regência, é a falta do sinal indicativo da crase – isso revela que o usuário não tem consciência de que, sob a forma do termo “ a ”, existe a presença de uma preposição “ a ”, exigida pela regência do termo anterior, e de um artigo definido.
Embora em outros exames, como o Enem, a falta de crase seja penalizada em acentuação, nesta avaliação esse desvio foi considerado no âmbito dos aspectos morfossintáticos.
Outro problema relacionado à regência verbal e à nominal, encontrado frequentemente nas questões, foi a ausência de preposição antes de pronome relativo, processo generalizado na modalidade oral da língua, em situações de registro informal.
Apesar da possibilidade de que essa alteração de regência se generalize no padrão escrito da Língua Portuguesa, como já está ocorrendo até em textos jornalísticos, o não emprego da preposição foi penalizado neste processo de avaliação.
Outro desvio muito frequente diz respeito aos processos de concordância verbal e de concordância nominal.
Quanto à concordância de número, observou -se ausência de marca com sujeito posposto ou uso indevido ( uso inadequado da marca de plural comandado pelo núcleo plural da locução adjetiva, apesar de o substantivo que funciona como núcleo do sintagma nominal estar no singular ).
Uma ocorrência generalizada foi a ausência de acento circunflexo na forma plural do presente do indicativo do verbo “ ter ”, que foi considerada como um desvio na concordância verbal e não na acentuação gráfica.
Quanto à concordância de gênero, vários casos foram observados, normalmente no âmbito de sintagmas nominais longos, em que o adjetivo está afastado do substantivo.
Quanto à questão da colocação pronominal, foram poucos os desvios observados.
Apesar de serem aspectos relacionados à oralidade, concluiu -se que, no registro escrito formal, a maioria dos participantes já incorporou regras como a não introdução da frase por um pronome oblíquo e a próclise na presença de um termo atrator.
Não se adotou, entretanto, o padrão excessivamente formal descrito pelas gramáticas normativas em relação à posição do pronome oblíquo em locuções verbais, já que esse uso está muito distante da prática cotidiana, até em textos mais formais.
Quanto aos aspectos vocabulares, alguns tipos de inadequação foram observados, principalmente na Questão 2, relacionada ao funk : expressões da oralidade, seleção vocabular incompatível com o contexto, gerando situações de falta de inteligibilidade ; falta de domínio de vocabulário mais abstrato e de maior complexidade, essencial ao desenvolvimento do texto de base dissertativa.
Ainda em função da temática desenvolvida na Questão 2, foi significativa a presença de palavras de baixa calão, por retratar um universo cultural específico.
Os textos que continham esse tipo de vocabulário foram considerados inadequados em função do registro formal exigido e foram avaliados com grau zero.
Várias marcas de oralidade foram identificadas, embora não em alta frequência : o uso do pronome relativo “ onde ” como relativo universal, falta de artigo definido antes de substantivo, repetição de palavras por falta de vocabulário, reduções como “ tá ” e “ pra ”, expressões informais. 3.3.2 Componente de Conhecimento Específico Na parte da prova relativa às questões discursivas do Componente de Conhecimento Específico ( Tabela 3.17 ), observa -se que a Média foi mais baixa do que para as questões discursivas do Componente de Formação Geral.
Enquanto no Componente de Formação Geral a Média para estudantes de Tecnologia em Gastronomia de todo o Brasil foi 51,8, na parte de Conhecimento Específico a Média foi 46,8.
A maior Média deste componente foi obtida pelos estudantes da região Sudeste ( 51,2 ), e a menor, pelos da região Norte ( 30,7 ).
Quanto à variabilidade das notas, o Desvio padrão de todo o Brasil foi 26,2.
O maior Desvio padrão foi encontrado na região Centro-Oeste ( 26,8 ), e o menor, na região Sudeste ( 25,2 ).
A nota Máxima ( 100,0 ) e a nota Mínima ( 0,0 ) foram obtidas por pelo menos um aluno em todas as regiões do Brasil, exceto a nota Máxima na região Norte ( 93,3 ).
A Mediana do Brasil como um todo foi 46,7, a mesma obtida nas regiões Sul e Centro-Oeste.
Na região Sudeste a Mediana foi 51,7, e nas regiões Norte e Nordeste, foi 28,3 e 38,3, respectivamente.
O Gráfico 3.10 representa a distribuição das notas nas questões discursivas no Componente de Conhecimento Específico.
A moda desta distribuição ocorre no intervalo ( 50 ; 60 ].
A análise de cada uma destas questões será feita a seguir. 3.3.2.1 Análise da Questão Discursiva 3 do Componente de Conhecimento Específico Na questão 3, cujos resultados aferidos encontram -se descritos na Tabela 3.18, a Média dos estudantes de todo o Brasil foi 53,8.
O desempenho dos estudantes nessa questão foi superior ao das questões 4 e 5.
A menor Média nessa questão foi obtida pelos alunos da região Norte ( 35,0 ), enquanto a maior Média foi obtida na região Sudeste ( 59,8 ).
Quanto à variabilidade das notas, o Desvio padrão de todo o Brasil foi 32,7.
O maior Desvio padrão foi obtido na região Norte ( 33,5 ), enquanto o menor foi obtido na região Sul ( 31,5 ).
A nota Máxima ( 100,0 ) e a nota Mínima ( 0,0 ) foram obtidas por pelo menos um aluno em todas as regiões do Brasil, sem exceção.
A Mediana do Brasil como um todo foi 55,0, a mesma obtida nas regiões Sudeste, Sul e Centro-Oeste, enquanto nas regiões Norte ( 30,0 ) e Nordeste ( 50,0 ), a Mediana foi menor.
O Gráfico 3.11 mostra a distribuição das notas na questão discursiva 3, do Componente de Conhecimento Específico da área de Tecnologia em Gastronomia.
Essa distribuição apresenta a moda no intervalo ( 70 ; 80 ], seguido de perto pelo intervalo ( 20, 30 ]. 3.3.2.2 Comentários sobre as respostas à Questão Discursiva 3 O perfil definido nas Diretrizes Curriculares Nacionais dos Cursos Superiores de Tecnologia em Gastronomia referencia profissionais com competências e habilidades para criar e elaborar pratos e cardápios, atuando direta e indiretamente na sua preparação, com conhecimento das ciências dos alimentos, sendo capazes de diferenciar e coordenar técnicas culinárias.
O item ‘ a ’ foi baseado no conhecimento que deveria ter sido adquirido pelo aluno durante o curso.
Pediu a composição de uma preparação básica muito usada na culinária, para dar sabor aos fundos e molhos na gastronomia.
O enunciado forneceu todas as informações necessárias para as respostas dos alunos, sendo bem claro e objetivo, e, ainda, sugeriu as opções dos ingredientes que podiam ser usados para o mirepoix, o que facilitou a elaboração da resposta do aluno.
Caberia a ele identificar a composição correta, usando as opções dos vegetais para a confecção do mirepoix, sugeridos no enunciado.
No entanto, mesmo assim, várias foram as respostas que incluíram outros vegetais e especiarias que não estavam citados no texto.
A questão foi de resposta simples, enquadrou -se perfeitamente no tempo de prova, e o espaço fornecido foi suficiente para os alunos responderem aos dois itens.
A questão pode ser classificada como de dificuldade média, quando pediu ao aluno, no item ‘ a ’, as proporções dos vegetais que deveriam ser usadas na composição do mirepoix.
As tentativas de resposta a essa solicitação elevaram bastante o número de erros observados na questão.
A maior diversidade de respostas em relação ao padrão estipulado ocorreu por erro de interpretação dos alunos ao que foi pedido no item ‘ b ’.
A pergunta se referiu a como os vegetais deveriam ser cortados em função do tempo de preparação, que foi colocado como menor do que o usual.
Os alunos fizeram associação com o tempo que se gastaria nos cortes dos vegetais, respondendo, então, que poderiam ser cortados em tamanho grande e de maneira irregular para não se perder tempo nesta fase, quando deveriam relacionar o tamanho dos vegetais com o tempo de cocção, devendo ser cortados em pequenos tamanhos para uma extração mais rápida do sabor, cor e aroma.
A correção do item ‘ a ’ foi bem objetiva na medida em que não dava muitas opções de resposta aos alunos.
Já o item ‘ b ’, que pedia os tipos de cortes com a respectiva justificativa, houve mais variação.
Vários alunos usaram linguagem apropriada, inclusive citando tipos de cortes que são usados pelo glossário gastronômico para responder ao item ‘ b ’.
Pelo número de respostas consideradas boas, observou -se que quase metade dos alunos que responderam à questão conseguiram entender o que era solicitado e responder de maneira correta aos itens ‘ a ’ e ‘ b ’.
Apesar de ser considerada uma questão de resposta fácil, foram encontrados números expressivos de respostas em branco e de notas zero.
Nas respostas fracas, os alunos só acertaram os ingredientes que deveriam ser usados no mirepoix, errando as proporções adequadas.
Nas respostas medianas, a predominância foi de acertos no item ‘ a ’, citando os vegetais e suas devidas proporções.
Poucos alunos erraram totalmente o item ‘ a ’, mas acertaram o item ‘ b ’, sugerindo os melhores cortes a serem usados na preparação.
As respostas consideradas boas foram em número bem significativo.
Nesses tipos de resposta os estudantes erraram apenas a justificativa pedida no item ‘ b ’, o que nos leva a supor que foi por erro de interpretação do enunciado.
Os principais erros cometidos pelos alunos nas respostas ao item ‘ a ’ foram as definições das proporções dos vegetais que deveriam ser usados no mirepoix.
A grande diversidade nas respostas foi quanto aos tipos de cortes que poderiam ser usados na preparação, o que significa que o assunto foi visto durante o curso.
No item ‘ b ’, apesar de a pergunta ser clara, muitos fizeram uma interpretação errada do que estava sendo pedido.
A questão não apresentou nenhuma similaridade com as questões objetivas da prova e estava muito bem inserida no que se espera do perfil do profissional de Gastronomia, que deve ser preparado para conhecer os fundamentos básicos das técnicas culinárias. 3.3.2.3 Análise da Questão Discursiva 4 do Componente de Conhecimento Específico A Tabela 3.19 contém as informações relativas à questão 4 do conjunto de questões do Componente de Conhecimento Específico.
O desempenho dos estudantes de todo o Brasil nesta questão foi mais baixo dentre as três questões discursivas desse componente.
A Média geral do Brasil foi 33,3, sendo a menor Média registrada na região Norte ( 14,8 ), e a maior, nas regiões Sudeste e Sul ( 37,2 ).
A nota Máxima ( 100,0 ) foi atingida por, pelo menos, um concluinte em cada uma das regiões do Brasil.
A Mediana e nota Mínima foram zero em todas as regiões, o que evidencia um alto percentual de alunos ( pelo menos 50 \% ) que obtiveram nota zero ou deixaram a questão em branco.
O Gráfico 3.12 representa a distribuição de notas da questão discursiva 4, do Componente de Conhecimento Específico.
Nessa distribuição, a moda é a categoria em branco, com mais do que 40 \% dos alunos nessa categoria. 3.3.2.4 Comentários sobre as respostas à Questão Discursiva 4 Apesar de estarem descritos no perfil do profissional de Gastronomia, o planejamento, o controle e a avaliação de custos na administração de restaurantes e estabelecimentos comerciais, são consideradas difíceis pelos alunos, já que exigem a resolução de situações que envolvem cálculos matemáticos.
Sem dúvida, foi a questão que apresentou o maior número de respostas em branco.
Observou -se, ainda, em algumas respostas, a tentativa dos alunos de resolverem os cálculos, porém, com muitos erros nas soluções apresentadas.
O enunciado da questão continha todas as informações necessárias para a sua resolução.
O espaço disponível foi suficiente para os dois itens.
E, a menos das dificuldades de cálculo, o tempo também deveria ter sido suficiente.
Vale ressaltar que o requisitado no item ‘ b ’ estava relacionado ao que se pedia no item ‘ a ’ e, dessa forma, respostas incorretas no item ‘ a ’ fizeram com que muitos alunos também perdessem a pontuação no item ‘ b ’.
O item ‘ a ’ solicitava que o aluno respondesse à questão efetuando três cálculos separadamente.
Muitos alunos não acertaram os três cálculos e, por isso, perderam todos os pontos relativos a este item, de acordo com o padrão de correção estabelecido junto à Comissão Assessora do Inep.
A questão teve a menor média entre as três questões discursivas, ainda que os valores tenham sido aproximados das demais.
Esta questão pode ser considerada a de mais difícil resolução da prova, como um todo, tanto entre as discursivas como entre as objetivas.
Foi, também, a questão que apresentou o maior número de notas 10,0 ( dez ) entre as questões discursivas de Conhecimentos Específicos, o que quer dizer que, se o aluno sabia fazer cálculos percentuais, não encontrava nenhuma dificuldade na sua resolução.
Essa questão apresentou número expressivo de notas zero.
As respostas consideradas como fracas, indicam que o aluno só acertou parcialmente o item ‘ b ’, errando o item ‘ a ’ integralmente.
As notas medianas obtidas pelos alunos referem -se, em sua maior parte, ao acerto total do item ‘ a ’ ou, em menor proporção, ao acerto integral do item ‘ b ’.
O número de provas com acerto nos itens ‘ a ’ e ‘ b ’ foi bastante expressivo, o que indica que o aluno que não tem dificuldades com cálculos percentuais não teve dificuldade em responder corretamente à questão.
Como mencionado anteriormente, os itens ‘ a ’ e ‘ b ’ estavam altamente correlacionados.
Os principais erros observados dos alunos que tentaram responder à questão foram os de cálculo.
Poucos alunos não sabiam o que era “ lucro líquido ” e ao que estava referido na pergunta.
Os alunos da Área de Gastronomia apresentaram bastante dificuldade na resolução de questões que envolviam cálculos matemáticos.
O perfil esperado do profissional de Gastronomia enfatiza que este deve ter competência para controlar e gerenciar custos, planejar, organizar, dirigir e controlar as atividades operacionais dos setores de produção e serviços em negócios de Gastronomia.
Logo, tendo como base os resultados deste exame, sugere -se que o conteúdo abordado nesta questão deva ser reforçado em todos os cursos, seja nas disciplinas de gestão, seja em outros componentes curriculares. 3.3.2.5 Análise da Questão Discursiva 5 do Componente de Conhecimento Específico A Tabela 3.20 contém as informações relativas à questão 5 do conjunto do Componente de Conhecimento Específico.
A nota Média dos estudantes de todo o Brasil foi 53,2.
A maior Média foi registrada na região Sudeste ( 56,7 ), enquanto a menor Média foi registrada na região Norte ( 42,2 ).
Quanto à variabilidade das notas, o Desvio padrão dos alunos do Brasil, como um todo, foi 31,5.
Enquanto o maior desvio foi encontrado na região Norte ( 33,7 ), o menor foi encontrado na região Sudeste ( 30,6 ).
A Mediana ( 60,0 ), e nota Máxima ( 100,0 ) e a nota Mínima ( 0,0 ) são as mesmas em todas as regiões, exceto a Mediana nas regiões Norte ( 40,0 ) e Nordeste ( 50,0 ).
O Gráfico 3.13 representa a distribuição das notas da questão discursiva 5 do Componente de Conhecimento Específico.
Nessa distribuição, a moda ocorre no intervalo ( 50 ; 60 ], seguido de perto pelo intervalo ( 70 ; 80 ]. 3.3.2.6 Comentários sobre as respostas à Questão Discursiva 5 Questão de conhecimentos básicos para todo profissional da Área de Gastronomia que trabalha, seja na manipulação, seja na preparação de alimentos.
Estes profissionais devem conhecer a importância dos princípios fundamentais da higiene pessoal, ambiental e de utensílios na manipulação de alimentos.
Devem estar cientes da sua responsabilidade no controle microbiológico e conhecer os métodos de conservação segundo os preceitos da Segurança Alimentar previsto nas Diretrizes Curriculares Nacionais dos Cursos Superiores de Tecnologia em Gastronomia.
O enunciado faz um resumo da cadeia produtiva de alimentos, destacando as etapas do recebimento e do processamento dos alimentos, sendo que, ao não especificar nenhum produto alimentício, aumentou, desta maneira, as possibilidades de resposta dos estudantes.
O padrão de resposta proposto pela Comissão Assessora da Área estava bastante diversificado, dando oportunidade para que pudessem, dentro das respostas esperadas, escolher apenas três componentes relacionados em cada um dos itens ‘ a ’ e ‘ b ’.
A questão possibilitava que os alunos usassem várias respostas para os itens ‘ a ’ e ‘ b ’.
Entretanto, observou -se, durante a correção, que vários alunos confundiram as etapas especificadas no enunciado : etapa de recebimento — item ‘ a ’ — e de processamento — item ‘ b ’.
Isso foi evidenciado pelas inúmeras respostas que se referiam às etapas de cocção e de armazenamento dos alimentos como se fizessem parte de um mesmo processo.
Foi uma questão adequada em relação ao tempo de prova, pois não exigia raciocínio muito elaborado do aluno.
O espaço também foi adequado, uma vez que não houve a necessidade de estender a resposta para além de seis a sete linhas para cada um dos itens da questão.
A temática envolvida nessa questão, a contaminação de alimentos, é adequada aos Cursos de Tecnologia de Gastronomia e pode ser abordada em diversas disciplinas da Área.
Apesar de o texto do enunciado fazer alusão a etapas diferenciadas, os alunos muitas vezes as confundiram, principalmente nas respostas ao item ‘ b ’.
A maioria respondeu de modo satisfatório à questão, principalmente ao item ‘ a ’, que pedia os cuidados a serem considerados na recepção dos alimentos.
Em algumas provas, ainda, observou -se uma deficiência na escrita dos estudantes, vários erros grosseiros de português e dificuldade de explicitar o pensamento de forma ordenada.
Em outras provas, observou -se que os alunos tiveram consciência dos cuidados a serem tomados, usando expressões corretas e pertinentes, \textbf{denotando} que este conteúdo foi bem trabalhado durante o curso.
Apesar de ser uma questão fácil de ser respondida, com várias possibilidades, observou -se uma alta prevalência de respostas em branco.
A média das notas ficou muito próxima da das demais questões, mesmo sendo a mais fácil dentre as discursivas de Conhecimentos Específicos.
Em relação ao item ‘ a ’, vários citaram a importância do prazo de validade, a integridade das embalagens, e os cuidados na temperatura dos produtos.
Muitos também fizeram alusão às características de higiene do transporte dos alimentos.
Em relação ao item ‘ b ’, as respostas mais frequentes foram em torno de se evitar o contato de alimentos crus e cozidos, evitando, assim, a contaminação \textbf{cruzada}.
Foi feita referência, também, ao uso de utensílios diferenciados na manipulação dos alimentos.
Os principais erros em relação ao item ‘ b ’, ocorreram por falta de clareza na separação das etapas, confundindo -se as de processamento com as de armazenamento.
Não se pode dizer que houve muita diversidade nas respostas que não tivessem presentes no padrão, que apresentava as várias opções que deveriam ser consideradas para ambos os itens.
Espera -se que a maioria dos alunos tenha conhecimento dos princípios básicos de higiene que devem ser adotados na manipulação de alimentos, cuidados que são importantes para se evitar a contaminação dos alimentos, seja a contaminação direta, causada pelos manipuladores de alimentos, como a contaminação \textbf{cruzada}, que foi objeto do item ‘ b ’ da Questão 5.
Por serem cuidados básicos, são sempre solicitados e destacados nas aulas práticas pelos professores de gastronomia.
Das três questões discursivas, sem dúvida, esta foi a questão que solicitou um cuidado maior na correção, principalmente por serem muito variadas as opções de resposta.
Na prova, não havia nenhuma outra questão que relacionasse os princípios básicos de higiene à Gastronomia.
Esta foi a única questão que abordou o assunto, apesar de sua importância na elaboração de preparações culinárias e de ser abordado em várias disciplinas do curso.

% matched lemmas: cruzar, denotar, indicado, visualizar
\end{textsample}
